\section*{Chapitre \ref{ch:rc:inflation-derivatives} --- Produits dérivés d'inflation}

\begin{solution}{exo:rc:inflation-derivatives:1}
Le ratio d'indice forward vaut $1.0328^{10} = 1.3809$~; le facteur d'actualisation réel est
$P_N(0,10)\times1.3809 = 0.6771\times1.3809 = 0.9349$.
\end{solution}

\begin{solution}{exo:rc:inflation-derivatives:2}
Le ratio de la première année $I(T_1)/I(0)$ a un dénominateur connu~: c'est un paiement zéro-coupon,
dont l'espérance sous la probabilité de paiement est exactement le ratio d'indice forward.
\end{solution}

\begin{solution}{exo:rc:inflation-derivatives:3}
Par 5~\% (le cap est actif)~; avec $[0~\%,2.5~\%]$, par 2.5~\%.
\end{solution}

\begin{solution}{exo:rc:inflation-derivatives:4}
L'espérance du ratio de l'année est le prix réel espéré $P_R(T,S)$ au
début de l'année, rapporté au prix nominal connu. L'incertitude sur les taux réels rend le prix
obligataire espéré (sous la probabilité utilisée) inférieur à son forward, au travers de la dérive $\mu(T)$ que
les taux réels stochastiques impliquent~: l'ajustement est négatif et croît avec $T$. Une corrélation positive, notée
$\rho$, elle abaisse la dérive du taux réel sous la probabilité nominale (le terme quanto), ce qui
élève le prix obligataire espéré et compense en partie cet effet.
\end{solution}

\begin{solution}{exo:rc:inflation-derivatives:5}
L'indice de mai~: de janvier à juin, le profil saisonnier américain ajoute environ un point au forward
par rapport à la trajectoire lisse, tandis qu'en décembre la somme du profil est nulle et le
forward revient sur la trajectoire lisse. Le fixing du swap de mai vaut donc davantage.
\end{solution}

\begin{solution}{exo:rc:inflation-derivatives:6}
La variation annuelle encadrée par le cap et le floor multiplie le montant accumulé jusque-là~: le paiement est
un produit d'options, dont la valeur dépend de la loi conjointe de l'ensemble des variations annuelles
et de la trajectoire, et non une somme d'options sur des années isolées.
\end{solution}

\begin{solution}{exo:rc:inflation-derivatives:7}
$-0.0182$ par unité lorsque $\sigma_I$ passe de 1.5~\% à 2.5~\%, $-0.0061$ lorsque $\sigma_R$ passe de 0.80~\%
à 0.90~\%~: la jambe est courte en volatilité (le cap à 5~\% vaut davantage pour le payeur à mesure que la volatilité
monte, plus que le floor à 0~\%).
\end{solution}

\begin{solution}{exo:rc:inflation-derivatives:8}
Le taux zéro-coupon capitalise un seul ratio d'indices forward~; le taux year-on-year moyenne
les espérances des ratios annuels, chacune décalée de son ajustement de convexité et pondérée par
les facteurs d'actualisation nominaux. Avec des taux réels stochastiques, les deux diffèrent~: 3.22~\% contre 3.28~\% à
dix ans ici.
\end{solution}

\begin{solution}{pb:rc:inflation-derivatives:1}
\textbf{1.} 0.8356 par simulation~; exactement $P_R(0,20) = 0.8345$ (l'écart provient de
l'erreur d'échantillonnage).
\textbf{2.} 0.7863.
\textbf{3.} 0.6371.
\textbf{4.} C'est l'indice avec un cap vendu à 5~\% et un floor acheté à 0~\% sur la variation de chaque année,
composé~; avec une inflation espérée proche de 3.3~\%, le cap est plus proche de la monnaie et coûte davantage que
ce que vaut le floor.
\textbf{5.} Une couverture non plafonnée surcouvre~: elle verse une inflation au-delà de 5~\% que les engagements n'encourent jamais,
donc le régime détient une position non couverte, longue en inflation et longue en volatilité, valant environ
$0.8356-0.7863 = 0.049$ par unité.
\textbf{6.} $-0.0182$ par unité.
\textbf{7.} $-0.0061$ par unité.
\textbf{8.} Détenir des engagements LPI et recevoir le LPI du swap laisse le régime couvert~; la
contrepartie qui paie le LPI est longue en volatilité (courte sur la jambe), puisque la jambe perd de la valeur à mesure que
la volatilité monte.
\textbf{9.} Afin que l'écart mesure l'effet du paramètre, et non le bruit de deux
simulations indépendantes.
\textbf{10.} Caps et floors year-on-year sur l'indice, et options sur indice là où elles se négocient~;
restent la corrélation et la dépendance au chemin.
\textbf{11.} 5~\%.
\textbf{12.} 14.2~\% sur l'année.
\textbf{13.} L'excédent de la couverture sur les engagements était un gain important, sans lien avec les obligations du
régime~: une position, et non une couverture, qui serait perdue si l'inflation s'inversait.
\textbf{14.} À partir de 2030, le RPI est calculé comme le CPIH, qui s'est situé en dessous~; le RPI forward
au-delà de 2030 baisse, et il en va de même pour la valeur des jambes LPI indexées sur le RPI.
\textbf{15.} Sans floor, la jambe ne comporte qu'un cap~: encore moins chère que la
jambe $[0~\%,5~\%]$, et le risque de déflation repose sur les membres du régime.
\textbf{16.} Les institutions ayant l'exposition inverse~: émetteurs d'actifs indexés sur le LPI (certaines
infrastructures, baux immobiliers), assureurs qui vendent des annuités, et dealers qui portent
la volatilité.
\textbf{17.} Il pilote la convexité et la loi conjointe des variations annuelles, et aucun instrument liquide
ne le cote~; son choix fait varier la valeur de la jambe.
\textbf{18.} La valeur des engagements plafonnés, le décalage avec les couvertures non plafonnées, la
sensibilité à la volatilité et à la corrélation, et un scénario de type 2022.
\textbf{19.} Résultat à retenir~: \emph{le swap LPI}~: la jambe non plafonnée vaut 0.8345 par unité (0.8356
simulé), la jambe LPI $[0~\%,5~\%]$ vaut 0.7863, et un point de volatilité supplémentaire de l'indice fait baisser la jambe LPI de
0.0182.
\textbf{20.} Elle transforme la promesse en un paiement truffé d'options, qui vaut moins que l'indice
et perd de la valeur à mesure que l'inflation devient plus volatile.
\end{solution}

\begin{solution}{iq:rc:inflation-derivatives:1}
La monnaie nominale est la devise domestique, une unité du panier est une devise étrangère, et
l'indice des prix est le taux de change. Les obligations réelles sont des obligations étrangères~; l'indice forward est
un taux de change forward $I(t)P_R/P_N$~; les taux réels vus depuis la probabilité nominale acquièrent une
dérive quanto. Les modèles multidevises deviennent des modèles d'inflation.
\iqlookfor{la correspondance et ses conséquences.}
\end{solution}

\begin{solution}{iq:rc:inflation-derivatives:2}
Chaque paiement est un ratio de deux valeurs futures de l'indice, versé à la date ultérieure~; son
espérance sous la probabilité de paiement est le ratio forward multiplié par un ajustement de convexité
qui dépend de la volatilité des taux réels et des corrélations, que la courbe zéro-coupon ne
contient pas.
\iqlookfor{ratio de valeurs futures et ajustement de convexité.}
\end{solution}

\begin{solution}{iq:rc:inflation-derivatives:3}
Les swaps et les obligations se réfèrent à l'indice d'un mois donné, non corrigé des variations saisonnières~; le forward
entre les points annuels doit porter le profil mensuel. Estimer les facteurs saisonniers à partir de
l'historique de l'indice (années complètes), conserver les points de swap annuels, et infléchir la trajectoire entre eux~;
vérifier la stabilité des facteurs dans le temps.
\iqlookfor{référence mensuelle et estimation historique.}
\end{solution}

\begin{solution}{iq:rc:inflation-derivatives:4}
Les hausses annuelles égales à la variation de l'indice, encadrées par le floor et le cap (au Royaume-Uni, $[0~\%,5~\%]$ ou
$[0~\%,2.5~\%]$), capitalisées. Les couvertures sont des options sur l'inflation à dépendance au chemin, peu liquides, avec
une exposition à la volatilité et à la corrélation~; les couvertures non plafonnées surcouvrent les années de forte inflation.
\iqlookfor{dépendance au chemin et leçon de 2022.}
\end{solution}

\begin{solution}{iq:rc:inflation-derivatives:5}
Observables~: les courbes nominales et réelles (swaps, linkers) et, moins liquides, les volatilités
d'indice issues des options year-on-year et zéro-coupon. Peu observables ou non~: le
retour à la moyenne et la volatilité du taux réel, et les corrélations entre taux nominaux, taux réels
et l'indice~; ils sont fixés à partir de l'historique ou implicites à partir de peu d'options.
\iqlookfor{quels paramètres sont posés et non calibrés.}
\end{solution}

\begin{solution}{iq:rc:inflation-derivatives:6}
La dérive de l'indice ($n-r^R-\frac12\sigma_I^2$), la dérive quanto du taux réel
($-\rho\sigma_R\sigma_I$), la construction de la courbe réelle à partir des swaps (capitalisation,
décalage), l'alignement de la grille temporelle sur les dates de paiement, et l'actualisation avec la courbe nominale~;
tester que l'indice simulé multiplié par le facteur d'actualisation nominal reproduit $P_R(0,T)$ à chaque
échéance.
\iqlookfor{dérives, probabilités et test de revalorisation directe.}
\end{solution}
